\section{Local coefficients and semisimple parameters}
\label{sec:coefficients}

The comparison in this section expresses the divisor of a Shahidi local
coefficient in terms of the semisimple local parameter.  It applies to
arbitrary generic inducing representations, including ramified
supercuspidals.  We use multiplicativity to reduce to supercuspidal data,
and then isolate the prescribed place in a global functional equation.
Only the product of local factors occurring in a local coefficient is
needed.  Lemma~\ref{lem:monodromy-ratio} allows this product to be compared
without identifying monodromy operators.

Let $H$ be a split adjoint group under consideration, or a standard Levi
subgroup of such a group, over the local function field $k$.  Fix a split
maximal torus $T$, a Borel subgroup $B=TU$, and a nondegenerate character
$\psi$ of $U(k)$.  Write $P=MR$ for a proper standard parabolic subgroup
of $H$.  All induction is normalized, and we write
\[
 I_P^H(\sigma_z)=\operatorname{Ind}_{P(k)}^{H(k)}(\sigma_z).
\]
Here $\sigma$ is an irreducible generic representation of $M(k)$ and
$z\in\mathfrak a_{M,\mathbb C}^*=X^*(M)\otimes_{\mathbb Z}\mathbb C$.
Whittaker characters and Weyl representatives are chosen compatibly.
Lemma~\ref{lem:whittaker} gives a single torus orbit of nondegenerate
characters for these groups, so such a choice is possible for every
generic $\sigma$.

Let $w_H$ and $w_M$ denote the longest Weyl elements of $H$ and $M$, and
put $w=w_Hw_M$.  Let $P'$ be the standard parabolic with Levi
$wMw^{-1}$.  The standard intertwining operator
\[
 A_H(z,\sigma;w):I_P^H(\sigma_z)
       \longrightarrow I_{P'}^H((w\sigma)_{wz})
\]
is defined by its usual unipotent integral and meromorphic continuation.
In particular, $A_H$ has no normalization by local
$L$-functions.  We use the convention
\begin{equation}
 \lambda_{\mathrm{source}}(z)
   =C_H(z,\sigma)\,\lambda_{\mathrm{target}}(wz)
                         \circ A_H(z,\sigma;w)
 \label{eq:coefficient-definition}
\end{equation}
for the local coefficient.  We also use this definition for a Weyl
element carrying the simple roots of its source Levi into the ambient
simple roots.  These are the intertwiners that occur in multiplicativity.

Let $\widehat P=\widehat M\widehat R$ be the corresponding standard
parabolic in $\widehat H$.  Decompose its unipotent Lie algebra under the
connected center of $\widehat M$:
\begin{equation}
 \widehat{\mathfrak n}
       :=\operatorname{Lie}(\widehat R)=\bigoplus_b V_b.
 \label{eq:central-weight-decomposition}
\end{equation}
The index $b$ runs over the weights that occur, $r_b$ is the representation
of $\widehat M$ on $V_b$, and $b(z)$ is the pairing with the central
cocharacter specified by the twist $z$.

\begin{proposition}[Comparison of local coefficients]
\label{prop:comparison}
For $H,P,M$, and $\sigma$ as above, let
$(\rho_\sigma,N_\sigma)$ be any Weil--Deligne pair in $\widehat M$ with
Weil part the semisimple local parameter of $\sigma$.  Then
\begin{equation}
 C_H(z,\sigma)\ \doteq\
 \prod_b
 \frac{L\bigl(1-b(z),r_b^\vee\circ(\rho_\sigma,N_\sigma)\bigr)}
      {L\bigl(b(z),r_b\circ(\rho_\sigma,N_\sigma)\bigr)}.
 \label{eq:coefficient-comparison}
\end{equation}
In particular, one may take $N_\sigma=0$.  The equality is an identity of
meromorphic functions of the unramified twists, up to a nonzero constant
times a monomial in twist coordinates.
\end{proposition}

We first explain how local coefficients behave under specialization and
induction in stages.  We then prove the rank-one and principal-series
cases of the proposition.  Globalization supplies the remaining
supercuspidal case.  The local-coefficient constructions and the crude
functional equation are those of Lomel\'i
\cite[\S\S1--2 and Theorem~4.3]{LomeliLS}; the globalization method is
that of Gan--Lomel\'i \cite[Theorem~1.1 and \S5]{GanLomeli}.
For the related construction of $\gamma$-factors from semisimple
parameters, see also \cite[Section~7]{GenestierLafforgue}.

\subsection{Whittaker lines in families}

For a smooth representation $V$ of $H(k)$, write $V_{U,\psi}$ for its
twisted coinvariants.  When $V=I_P^H(\sigma_z)$ with $\sigma$ generic,
these coinvariants are one-dimensional.  The scalar of the standard
intertwining operator on this line is the inverse local coefficient.
To use this description at a reducibility point, one must choose
generators that remain nonzero after specialization.

Choose a finite cover of the torus of unramified twists on which the
characters in use are algebraic, and let $\mathcal A$ be its Laurent
polynomial coordinate ring.  The family $\sigma_z$ is a smooth
$\mathcal A[M(k)]$-module, and its induced family is a smooth
$\mathcal A[H(k)]$-module.  Meromorphic operators are obtained by
extending scalars to the fraction field of $\mathcal A$.

\begin{lemma}[Whittaker lines and specialization]
\label{lem:whittaker-family}
The twisted coinvariants of the induced family $I_P^H(\sigma_z)$ form a
free rank-one $\mathcal A$-module, and this description commutes with
specialization of $z$.  The Jacquet-integral Whittaker functionals give a
generator of its dual whose specialization is nonzero at every twist.
Consequently, the inverse scalar induced by $A_H$ on these lines is
$C_H(z,\sigma)$ up to a unit of $\mathcal A$.
\end{lemma}

\begin{proof}
The twisted Jacquet functor is exact over $\mathcal A$.  Indeed, the
unipotent group $U(k)$ is an increasing union of compact open subgroups.
On each such subgroup, twisted averaging is an idempotent, since
$\mathcal A$ is a complex algebra.  Taking the resulting filtered
colimit proves exactness.  The same description shows that twisted
coinvariants commute with scalar extension, including specialization.

Apply this functor to the Bruhat filtration of parabolic induction.
For a representative $x$ of $W_M\backslash W_H$, the corresponding
subquotient consists of sections with compact support modulo $P(k)$
on the cell $P(k)\dot xU(k)$.  As a $U(k)$-representation, it is compact
induction from
\[
 U(k)\cap\dot x^{-1}P(k)\dot x,
\]
with the fiber action obtained from the inducing representation.
Its twisted coinvariants are therefore computed on the fiber.  All
unramified twists and modulus characters are trivial on the unipotent
groups in this calculation, so the calculation is unchanged over
$\mathcal A$.

Every cell other than the open one has a simple root subgroup of $U$
whose conjugate lies in $R$.  It acts trivially on the fiber and
nontrivially through $\psi$, so that cell contributes no twisted
coinvariants.  For completeness, take $x$ minimal on the left modulo
$W_M$.  If no simple root is sent to a positive root outside $M$, every
$x\alpha$ for $\alpha$ simple is either negative or a positive Levi
root. Thus the coefficients outside the Levi simple system of every
$x\beta$, $\beta$ positive, are nonpositive. In particular, $x\beta$
cannot be a positive root outside $M$. It follows that $x^{-1}$ sends
all positive roots outside $M$ to negative roots, which characterizes
the representative of the open cell.

On the open cell, the remaining coinvariant calculation is precisely
the Whittaker coinvariant space of $\sigma$, with the character
transported by the chosen representative.  This space is
one-dimensional by Whittaker uniqueness.  Hence the induced family's
coinvariants are its tensor product with $\mathcal A$, as asserted.
This is the usual open-cell proof of Whittaker heredity, now carried
out over the coefficient ring.

The open-cell submodule therefore induces an isomorphism on twisted
coinvariants.  Its unipotent integral, applied to a fixed Whittaker
functional of $\sigma$, extends uniquely through this isomorphism to
an $\mathcal A$-linear Whittaker functional on the entire induced
family.  In a convergence chamber this functional and the usual
Jacquet integral agree on the open-cell submodule, and hence agree
everywhere.  Meromorphic continuation consequently identifies the
Jacquet functional with this algebraic functional.  In particular,
its evaluations on algebraic sections are Laurent polynomials, as
also recalled in \cite[\S1.2]{LomeliLS}.

To see that specialization is nonzero, choose a vector on which the
fixed Whittaker functional of $\sigma$ is $1$ and a compactly supported
function in the open-cell unipotent coordinates whose $\psi$-weighted
integral is $1$.  The resulting algebraic section has Whittaker value
$1$ at every twist.  Thus the regular functional just constructed
never specializes to zero and generates the dual coinvariant line.

Equation~\eqref{eq:coefficient-definition} now identifies the local
coefficient with the inverse scalar on these free lines.  Two
generators of a free rank-one Laurent polynomial module differ by a
unit, hence by a nonzero constant times a monomial.  Haar measures and
compatible Weyl representatives change the formula only by such
units.  Torus conjugation of the Whittaker character also identifies
these lines over $\mathcal A$ and commutes with the intertwining
calculation, with the same consequence.
\end{proof}

\subsection{Elementary intertwiners and reduction of the inducing data}

We record the root decomposition that governs both sides of
\eqref{eq:coefficient-comparison}.  For a Weyl element $v$, set
\[
 \operatorname{Inv}(v)
   =\{\beta\in\Phi_H^+:v\beta\in-\Phi_H^+\}.
\]
For a standard Levi $L$, write $\Delta_L$ for its simple roots.

\begin{lemma}[Elementary factorization]
\label{lem:elementary-moves}
Let $L$ be a standard Levi of $H$, and suppose $v\Delta_L\subseteq
\Delta_H$.  There is a length-additive factorization
\[
 v=v_t\cdots v_1
\]
with the following properties.  Put $v_{<j}=v_{j-1}\cdots v_1$ and
$L_j=v_{<j}Lv_{<j}^{-1}$.  Then $L_j$ is standard, it is a maximal
proper Levi in a standard Levi subgroup $H_j$, and
\[
 \Delta_{H_j}=\Delta_{L_j}\cup\{\alpha_j\},\qquad
 v_j=w_{H_j}w_{L_j}.
\]
Moreover,
\begin{equation}
 \operatorname{Inv}(v)
   =\bigsqcup_{j=1}^t v_{<j}^{-1}\operatorname{Inv}(v_j).
 \label{eq:inversion-partition}
\end{equation}
The analogous partition holds for coroots.
\end{lemma}

\begin{proof}
If $v\ne1$, choose a simple root $\alpha$ with $v\alpha<0$.
Such a root is outside $\Delta_L$, since $v\Delta_L\subseteq\Delta_H$.
Let $H_1$ be the standard Levi obtained by adjoining $\alpha$ to
$\Delta_L$, and set $v_1=w_{H_1}w_L$.

Every positive root of $H_1$ outside $L$ has a strictly positive
$\alpha$-coefficient.  Its image under $v$ is a positive linear
combination of $v\alpha$ and roots in $v\Delta_L$.  The negative root
$v\alpha$ has a strictly negative coefficient outside $v\Delta_L$:
otherwise it would belong to the span of $v\Delta_L$, contradicting
$\alpha\notin\operatorname{span}(\Delta_L)$.  This coefficient cannot
be canceled by the other summands.  The image root is therefore
negative.  It follows that
\[
 \operatorname{Inv}(v_1)
   =\Phi_{H_1}^+\setminus\Phi_L^+
   \subseteq\operatorname{Inv}(v).
\]
Consequently $v_1$ is a right factor with additive length.

The element $v_1$ carries $\Delta_L$ into the simple roots of $H_1$.
Replace $L$ by $v_1Lv_1^{-1}$ and $v$ by $vv_1^{-1}$, which still
carries the new Levi simple system into $\Delta_H$.  The length has
strictly decreased, so iteration terminates.  The usual inversion-set
identity for a product with additive lengths gives
\eqref{eq:inversion-partition}.  Applying the same argument to the dual
root system gives its coroot version.
\end{proof}

The standard intertwining integrals factor according to
Lemma~\ref{lem:elementary-moves}.  First one uses Fubini's theorem in
a convergence chamber, and then meromorphic continuation.  At each
step, normalized induction in stages identifies the elementary
operator with the maximal-parabolic operator in $H_j$.  Taking
Whittaker scalars and using Lemma~\ref{lem:whittaker-family} therefore
gives multiplicativity of local coefficients, up to the allowed
units.  This is the multiplicativity of
\cite[Proposition~2.3]{LomeliLS}.

We explain more precisely how it reduces the inducing representation.
Suppose for the moment that $P$ is maximal in $H$, with omitted simple
root $\alpha$.  Let $\widetilde\alpha\in\mathfrak a_{M,\mathbb R}^*$
be proportional to the character giving the determinant on
$\operatorname{Lie}(R)$ and normalized by
\[
 \langle\widetilde\alpha,\alpha^\vee\rangle=1.
\]
Its pairing with every simple coroot of $M$ is zero.  The decomposition
\eqref{eq:central-weight-decomposition} is then conventionally written
\[
 \widehat{\mathfrak n}=\bigoplus_{i\geq1}V_i,\qquad
 V_i=\bigoplus_{\langle\widetilde\alpha,\beta^\vee\rangle=i}
                       \widehat{\mathfrak h}_{\beta^\vee},
\]
where $\widehat{\mathfrak h}_{\beta^\vee}$ is the corresponding root
space in $\operatorname{Lie}(\widehat H)$ and $i$ ranges over the
positive integers that occur.  Along
$z=s\widetilde\alpha$, the arguments in
\eqref{eq:coefficient-comparison} are $is$ and $1-is$.

Every other twist direction is a sum of this relative direction and a
direction in $X^*(H)\otimes\mathbb C$.  A twist from $H$ twists both
sides of the intertwiner and changes its Whittaker scalar by at most
a unit.  On the dual side it acts through $Z(\widehat H)^\circ$ and
acts trivially on $\widehat{\mathfrak n}$.  Thus it suffices to prove
the maximal-parabolic comparison on the relative line.

By the subrepresentation theorem, $\sigma$ embeds in normalized
induction in $M$ from a supercuspidal representation $\sigma_0$ of a
standard Levi $M_0\subseteq M$.  Exactness and heredity of Whittaker
coinvariants show that $\sigma_0$ is generic.  They also show that the
embedding induces an isomorphism on Whittaker lines, since both lines
are one-dimensional.  The same holds after every unramified twist.

The element $w=w_Hw_M$ carries the simple roots of $M_0$ into those of
$H$.  In the induced realization from $M_0$, its intertwining integral
still uses exactly the roots outside $M$ turned negative by $w$:
$w$ preserves the positive system of $M$ in its target Levi.  Thus
induction in stages identifies the intertwiner on $\sigma$ with the
restriction of this intertwiner on the smaller inducing data.  The
isomorphisms of Whittaker lines just established identify their local
coefficients, up to units.

Apply Lemma~\ref{lem:elementary-moves} with $L=M_0$, and write
$M_j=w_{<j}M_0w_{<j}^{-1}$ for the successive Levi subgroups.  At step $j$,
the inducing representation is the transport of $\sigma_0$ to $M_j$,
and the twist is $w_{<j}(s\widetilde\alpha)$.  Its relative coordinate
in $H_j$ is $c_js$, where
\begin{equation}
 c_j=\langle w_{<j}\widetilde\alpha,\alpha_j^\vee\rangle>0.
 \label{eq:relative-slope}
\end{equation}
Indeed, $w_{<j}^{-1}\alpha_j^\vee$ belongs to the original inversion
set of coroots, all of which are positive and outside $M$.  Pairing
with $\widetilde\alpha$ is strictly positive on this set.  Each
elementary coefficient consequently has a genuinely varying argument;
the reduction does not restrict a meromorphic coefficient to an
identically singular locus.

The same factorization decomposes the dual representations.  Transport
the nilradical Lie algebra for $M_j\subset H_j$ back by $w_{<j}^{-1}$.
It is invariant under $\widehat M_0$, because the original algebra is
invariant under $\widehat M_j$.  The coroot partition
\eqref{eq:inversion-partition} identifies their direct sum with
$\widehat{\mathfrak n}$ restricted to $\widehat M_0$.  A step summand
of relative degree $i_j$ has twist exponent $c_ji_js$.  To see this,
decompose $w_{<j}\widetilde\alpha$ into
$c_j\widetilde\alpha_j$ and a character direction from $H_j$; the
latter pairs to zero with every coroot in $H_j$.

Compatibility of semisimple parameters with normalized induction
identifies $\rho_\sigma$ with the parameter obtained from
$\rho_{\sigma_0}$.  Taking monodromy zero, the preceding decomposition
therefore makes the right side of
\eqref{eq:coefficient-comparison} the product of the right sides for
the elementary maximal parabolics.  We have proved that comparison
for maximal parabolics with supercuspidal inducing data implies
comparison for maximal parabolics with arbitrary generic data.

\subsection{The principal-series comparison}

The rank-one case supplies the local comparisons needed away from the
place prescribed in globalization.  If a maximal proper Levi $M$ is a
torus, the derived root system of $H$ has rank one.  The local
coefficient is the Tate gamma factor for the character obtained by
pulling $\sigma$ back along the coroot.  Tate's formula and local class
field theory give
\begin{equation}
 C_H(s\widetilde\alpha,\sigma)
   \ \doteq\
   \frac{L(1-s,r_1^\vee\circ\rho_\sigma)}
        {L(s,r_1\circ\rho_\sigma)}.
 \label{eq:rank-one-comparison}
\end{equation}
This calculation permits an arbitrary smooth inducing character;
see \cite[\S1.3]{LomeliLS}.  It also applies to the isogeny types
occurring here.  Pulling the integral back through the root
homomorphism from $\operatorname{SL}_2$ identifies the root groups
and the unipotent integration, and the inducing character pulls back
through the stated coroot.  Extra central tori have no effect on this
calculation.

An irreducible generic constituent of a principal series can be
embedded in a principal series, with a possibly Weyl-conjugate
inducing character, by the subrepresentation and supercuspidal
support theorems.  The preceding elementary factorization with
$M_0=T$ reduces its comparison to
\eqref{eq:rank-one-comparison}.  Thus the maximal-parabolic case of
\eqref{eq:coefficient-comparison} is already proved whenever the
inducing representation is a generic principal-series constituent.
In particular it is available for ramified principal series; an
Iwahori-fixed hypothesis is not needed.

\subsection{Globalization of supercuspidal data}

We now prove the remaining maximal-parabolic comparison.  Let
$\sigma$ be a generic supercuspidal representation of $M(k)$.  We
may first twist it so that its central character has finite order.
Indeed, $Z(M)$ is a split torus, a smooth character of its maximal
compact subgroup has finite image, and restriction
\[
 X^*(M)\longrightarrow X^*(Z(M))
\]
has finite cokernel.  Unramified twists can therefore adjust the
finitely many uniformizer values to roots of unity.  Twisting back
shifts the relative parameter $s$; the remaining twist direction is
from $H$ and has the behavior already described.  Twist compatibility
of semisimple parameters gives the identical shift on the proposed
Galois expression.

Write $k=k_0((t))$, and put $K=k_0(t)$.  Let $v_0$ be the rational
place $t=0$, so $K_{v_0}=k$.  Use the constant split models of $H$ and
$M$ over $K$.  The finite-order central character just obtained
extends to a finite-order automorphic character of $Z(M)(\mathbb A_K)$.
Here is an explicit extension, to check the central-character
hypothesis of globalization.

For one split central factor, the idele-class group of $K$ fits into
the split exact sequence
\begin{equation}
 1\longrightarrow
    \left(\prod_x\mathcal O_x^\times\right)/k_0^\times
   \longrightarrow K^\times\backslash\mathbb A_K^\times
   \xrightarrow{\deg}\mathbb Z\longrightarrow0.
 \label{eq:rational-idele-class-group}
\end{equation}
The degree-zero description follows from the triviality of
$\operatorname{Pic}^0(\mathbb P^1)$, and the uniformizer idele at
$v_0$ splits the degree map.  Prescribe the given unit character at
$v_0$, cancel its restriction to $k_0^\times$ using a character of
the residue field at infinity, and take trivial characters at the
remaining unit groups.  This defines a finite-order character of
the left group in \eqref{eq:rational-idele-class-group}.  Give the
degree generator the prescribed finite-order value of the local
character on $t$.  The resulting automorphic character has the
desired restriction at $v_0$.  Repeat for every split central factor.

Choose a nontrivial global additive character and the resulting
generic character of the Borel unipotent of $M$.  Its local character
at $v_0$ is in the torus orbit for which $\sigma$ is generic.  The
globalization theorem of Gan--Lomel\'i
\cite[Theorem~1.1]{GanLomeli} now produces a globally generic
cuspidal representation $\Pi$ of $M(\mathbb A_K)$ such that
\[
 \Pi_{v_0}\simeq\sigma,
 \qquad
 \Pi_x\text{ is a principal-series constituent for every }x\ne v_0,
\]
and whose central character is the finite-order character just
constructed.  Their theorem preserves the prescribed unipotent
period; for the Borel unipotent and a nondegenerate character this
is precisely global genericity.  It imposes no restriction on the
depth of $\sigma$ or on the positive characteristic.

Choose a global parameter $\Sigma$ occurring for $\Pi$, and enlarge
a finite set $S$ of places containing $v_0$ so that normalized
Satake matching and all unramified conditions for the crude
functional equation hold outside $S$.  For the maximal parabolic
under consideration, Lomel\'i's crude functional equation is
\begin{equation}
 \prod_i
 \frac{L^S(is,\Pi,r_i)}
      {L^S(1-is,\Pi,r_i^\vee)}
     \ \doteq\ \prod_{x\in S}
                 C_{H_x}(s\widetilde\alpha,\Pi_x).
 \label{eq:crude-functional-equation}
\end{equation}
In \cite[Theorem~4.3]{LomeliLS}, the denominator is written using
the contragredient $\widetilde\Pi$ and $r_i$.  At an unramified
place, contragredience inverts the Satake class up to the Levi Weyl
group, so those Euler factors are exactly the factors for $\Pi$
and $r_i^\vee$ displayed here.

The same partial Euler products are the partial $L$-functions of
$r_i\circ\Sigma$.  Write $\Phi_x$ for the local Weil--Deligne
parameter obtained from $\Sigma$ at $x$.  The functional equation
for a lisse sheaf on a curve, with ramified factors defined by
inertia invariants, gives
\begin{equation}
 \prod_i
 \frac{L^S(is,\Pi,r_i)}
      {L^S(1-is,\Pi,r_i^\vee)}
 \ \doteq\
 \prod_{x\in S}\prod_i
   \frac{L(1-is,r_i^\vee\circ\Phi_x)}
        {L(is,r_i\circ\Phi_x)}.
 \label{eq:galois-functional-equation}
\end{equation}
This is the usual curve functional equation
\cite[\S\S9--10]{DeligneConstants}; it follows from the trace formula
and duality for middle extensions, and does not require purity.
To check the direction of the local ratio, write the complete
functional equation as $L(s,V)=\varepsilon(s,V)L(1-s,V^\vee)$.
Removing the Euler factors in $S$ puts
$L_x(1-s,V^\vee)/L_x(s,V)$ on the right.  The global epsilon factor
is a nonzero constant times a monomial.  After the substitutions
$s\mapsto is$, its contribution has exactly the form suppressed
by $\doteq$.

The $\ell$-adic functional equation is an identity of rational
functions and is transported through the fixed coefficient
identification.  Its partial $L$-functions agree with the
automorphic ones because their Euler factors agree outside $S$.
The ramified factors are those of $\Phi_x$: the invariant stalk of
the lisse sheaf is the inertia-invariant kernel of its monodromy,
and Frobenius semisimplification does not change its determinant.

We can now isolate $v_0$.  For each $x\in S\setminus\{v_0\}$, the
representation $\Pi_x$ is generic and is a principal-series
constituent.  The principal-series comparison identifies its
factor in \eqref{eq:crude-functional-equation} with its factor in
\eqref{eq:galois-functional-equation}, up to a unit.  Local-global
compatibility, together with Lemma~\ref{lem:wd-semisimplification},
identifies the Weil part of the Frobenius-semisimplified pair $\Phi_x$
with $\rho_{\Pi_x}$.  Lemma~\ref{lem:monodromy-ratio} then permits
us to discard its monodromy in comparing the local ratios.
Cancellation leaves
\[
 C_H(s\widetilde\alpha,\sigma)
   \ \doteq\
   \prod_i
   \frac{L(1-is,r_i^\vee\circ\Phi_{v_0})}
        {L(is,r_i\circ\Phi_{v_0})}.
\]
At $v_0$ the same compatibility and
Lemma~\ref{lem:monodromy-ratio} replace $\Phi_{v_0}$ by
$(\rho_\sigma,0)$, or by any specified completion
$(\rho_\sigma,N_\sigma)$.  This proves
\eqref{eq:coefficient-comparison} for maximal parabolics with
supercuspidal data.  The reduction above proves it for every
generic inducing representation of a maximal Levi.

\subsection{Completion of the comparison for an arbitrary parabolic}

\begin{proof}[Proof of Proposition~\ref{prop:comparison}]
The maximal-parabolic case has just been proved.  For general $P=MR$,
apply Lemma~\ref{lem:elementary-moves} to $w=w_Hw_M$, now with
$L=M$, and put $M_j=w_{<j}Mw_{<j}^{-1}$.  At step $j$, the inducing data are transported from
$\sigma$ to $M_j$, and the twist is $w_{<j}z$.  Its relative
coordinate is
\[
 s_j(z)=\langle w_{<j}z,\alpha_j^\vee\rangle.
\]
This linear function is strictly positive on the positive chamber
for $P$, because its transported coroot lies outside $M$ in the
inversion partition.  In particular it is nonconstant.
The maximal-parabolic comparison is available for each step and
for every such twist, with directions from $H_j$ contributing
only units.  Multiplicativity gives the product of these step
formulas for $C_H(z,\sigma)$.

Use monodromy zero in these formulas.  Transporting the step
nilradicals back to $\widehat M$, the inversion partition gives
their direct sum as $\widehat{\mathfrak n}$.  On a step summand
of degree $i_j$, the twist exponent is $i_js_j(z)$, which is
exactly $b(z)$ for its connected-central weight in
\eqref{eq:central-weight-decomposition}.  Multiplicativity of
Artin $L$-factors for direct sums identifies the product of all
step expressions with the right side of
\eqref{eq:coefficient-comparison} for $N_\sigma=0$.
Finally Lemma~\ref{lem:monodromy-ratio}, applied to each $r_b$,
allows any Weil--Deligne completion with the same Weil part.
\end{proof}
